问题四的检验在 Q2/Q3 已有审计之上，增加\textbf{价格因果性}与\textbf{参数校准一致性}两类硬约束。

\paragraph{信息因果性。}日初与每次日内预测的训练样本日期必须严格早于决策时点；对未来真实价格、负荷或光伏作扰动测试时，已锁定的 $q$、$g^0$ 与已执行轨迹不得改变。价格预测日志须可逐日复算候选 MAE 与在线选择结果。

\paragraph{三元路径一致性。}每个代表情景的 $(L,P,p)$ 必须来自同一历史日，报告相关性与代表路径覆盖范围，但不将统计相关叙述为因果机制。

\paragraph{物理与结算审计。}对 Q4-2 抽查 $x\le q$、能量平衡、SOC 递推、功率与弃光边界、跨日 SOC 与同时充放电；对 Q4-3 另查 6:00/12:00/18:00 后 $g^F$ 只影响之后时段、结算始终相对 $g^0$。实现成本必须分别按式\eqref{eq:q4-2-actual-cost} 与式\eqref{eq:q4-3-actual-cost} 重算，并与 \texttt{q4\_2\_cost\_ledger.csv}（及 Q4-3 台账）一致。

\paragraph{RQ4-C1 与 $K$ 复核。}每个 $\alpha$ 候选的评分必须采用与部署相同的“联合情景计划—锁定额度—逐十分钟 MPC”链条，并连续传递候选自身 SOC；禁止用整日真实路径的全视域 LP 选参。在相同闭环协议下比较 $K=4,8,12$ 的实际成本、紧急购电量与耗时，说明保留 $K=8$ 的判据。

\paragraph{Q4-2 当前状态。}正式全年运行已完成 RQ4-C1 重标定，且校准未调用整日真实轨迹的全视域调度：365 日物理审计全部通过，跨日 SOC 连续；两个代表性日期的信息集扰动抽查与 $\alpha$ 的未来扰动不变性测试均通过。该抽查用于发现信息泄漏，不替代逐日历史截止日、锁定额度和物理审计。$K=8$ 在 25 个校准窗口内相对 $K=12$ 的验证成本仅高 $0.243\%$，平均耗时由 $6.7$ s 降至 $5.6$ s（约 $16.4\%$）；在本文采用的 $1\%$ 成本容差下，取其为在线可用的稳定折中，而非成本最优。正式工作簿 \texttt{result4-2.xlsx} 已与对应成本台账一致导出。

\paragraph{Q4-3 当前状态。}Q4-3 全年链路已完成：365 日物理审计全部通过，跨日 SOC 连续；两个代表性日期的信息集扰动抽查通过，逐时段成本重算与日台账一致。正式工作簿 \texttt{result4-3.xlsx} 已按 334 日题设区间导出；与 Q3 固定价主结论的比较在正文中仅作\textbf{价格机制对照}，不作同口径优劣排序。
